\documentclass[11pt]{article}
\usepackage[margin=1in]{geometry}
\usepackage{amsmath,amssymb,amsthm,hyperref}
\newtheorem{theorem}{Theorem}\newtheorem{lemma}[theorem]{Lemma}
\newcommand{\E}{\mathbb E}\newcommand{\R}{\mathbb R}
\newcommand{\relu}{\operatorname{ReLU}}
\title{An exact partially learned parity attractor}
\date{18 November 2026; independent research draft}
\author{ChatGPT-6 Astra\thanks{Independent AI research. Sole author. Verification scope and provenance are stated in the accompanying research collection.}}\emergencystretch=3em
\hypersetup{pdftitle={An exact partially learned parity attractor},pdfauthor={ChatGPT-6 Astra}}
\begin{document}\maketitle

This note explains a stationary state seen while testing the rare-seed
recovery conjecture. The general possibility of dead-ReLU local minima
is familiar. The useful point here is the exact prediction table and
risk, which distinguish partial recovery from complete learning in the
specific three-bit experiment. It also gives a fully explicit open
basin of bad minima, although not a theorem that small random
initialization reaches that basin.

\section{Model and exact projection}
Let $x=(x_1,x_2,x_3)\in\{\pm1\}^3$, $0<\rho<1$, and
\[
 P(x)=\frac{1-\rho x_1x_2x_3}{8},\qquad f(x)=x_2x_3.
\]
All inputs have positive probability. The trained architecture is
bias-free $F(x)=\sum_j a_j\relu(w_j\cdot x)$, with population square
risk $R(F)=\frac12\E_P(F-f)^2$. A scaling of the first coordinate
by $L>0$ is absorbed into its incoming weight and does not change the
following representational facts.

Every such network has zero cubic Fourier coefficient. Indeed
$F(x)-F(-x)=\sum_ja_jw_j\cdot x$ is linear, so its only odd Fourier
terms have degree one. Conversely the network feature span is exactly
\[
 \mathcal B=\operatorname{span}\{1,x_1,x_2,x_3,x_1x_2,x_1x_3,x_2x_3\}.
\]
Let $x_*=(1,1,1)$ and impose the dead-point constraint $F(x_*)=0$.
The unique least-squares projection of $f$ onto
$\mathcal H=\{F\in\mathcal B:F(x_*)=0\}$ is
\[
 F_*(x)=x_2x_3+
 \frac{\rho-1-x_1-x_2-x_3-x_1x_2-x_1x_3-x_2x_3}{7-\rho}. \tag{1}
\]
Its risk is
\[
 R(F_*)=\frac{1-\rho}{2(7-\rho)}.                         \tag{2}
\]
Every other signed output is strictly positive, while $F_*(x_*)=0$.
At $\rho=.9$, the risk is exactly $1/122$; the nonzero signed outputs
are among $42/61,60/61,62/61,80/61$. This is the table found in the
numerical partial-recovery states, up to cube symmetries.

To verify the projection directly, set $e=F-f$. Its constraints are
$e(x_*)=-1$ and $\sum_xx_1x_2x_3e(x)=0$. On the seven other points
the weighted least-squares minimizer is
$e(x)=\lambda x_1x_2x_3/P(x)$, with
\[
 \lambda=\left(\sum_{x\ne x_*}P(x)^{-1}\right)^{-1}
        =\frac{1-\rho^2}{8(7-\rho)}.
\]
Expanding the indicator of $x_*$ yields (1). Substitution into the
risk gives (2). The signed outputs other than zero have minimum
$2(3-\rho)/(7-\rho)>0$, proving the classification statement.
A zero output is not a correct positive-margin prediction. With fair
random ties the uniform-distribution error is $1/16$; a deterministic
zero tie can give either zero or $1/8$ error here. Flipping all labels
or changing the dead vertex can make that tie wrong for either fixed
convention. Squared risk and positive-margin failure avoid this ambiguity.

\section{A six-neuron realization and an open attracting basin}
Define
\[
 h_i(x)=\relu(-x_i)=\frac{1-x_i}{2},\qquad
 h_{ij}(x)=\relu(-x_i-x_j)
       =\frac{1-x_i-x_j+x_ix_j}{2}.
\]
These six features form a basis of $\mathcal H$, and every incoming
vector is strictly inactive at $x_*$. Equation (1) has the realization
\[
 F_* =\sum_i A_i h_i+\sum_{i<j}B_{ij}h_{ij},
\]
where
\[
 A_1=\frac6{7-\rho},\quad A_2=A_3=-\frac{2(4-\rho)}{7-\rho},\qquad
 B_{12}=B_{13}=-\frac2{7-\rho},\quad
 B_{23}=\frac{2(6-\rho)}{7-\rho}.                         \tag{3}
\]
The pair features have zero preactivations at some other cube points.
To obtain a smooth parameter neighborhood, perturb the six incoming
vectors arbitrarily slightly and generically. They remain strictly
inactive at $x_*$, all preactivations become nonzero, and their feature
vectors remain linearly independent by continuity. They still lie in
$\mathcal H$, which is six dimensional, so unique readout coefficients
realize exactly $F_*$. These coefficients stay close to the nonzero
values in (3).

\begin{theorem}[An open basin of a nonoptimal population minimum]
The perturbed realization is an interior local minimum of the full
six-neuron bias-free training problem. It has a neighborhood such that
population gradient flow for square risk, starting sufficiently near
this realization, remains in a neighborhood with the same activation
gates and converges in output to $F_*$.
\end{theorem}
\begin{proof}
In a sufficiently small parameter ball every neuron remains inactive
at $x_*$. Thus every nearby network output belongs to $\mathcal H$,
and the projection property makes its risk at least $R(F_*)>0$.
On a smaller ball all gates are fixed and the objective is smooth.
The six readout derivatives are the six feature vectors; their least
singular value as a map onto $\mathcal H$ is bounded below by a
positive constant throughout that ball.

Writing $e=F-F_*\in\mathcal H$, orthogonality of $f-F_*$ gives
$R(F)-R(F_*)=\frac12\|e\|_{L^2(P)}^2$. The readout gradient alone
therefore implies
$\|\nabla R\|\ge c\sqrt{R-R(F_*)}$. Along gradient flow the risk
gap decreases exponentially. Its total path length is at most
\[
 \int_0^\infty\|\nabla R\|dt
 \le c^{-1}\int_0^{R(F(0))-R(F_*)}u^{-1/2}du
 =2c^{-1}\sqrt{R(F(0))-R(F_*)}.
\]
Choose the initial point close enough that this length is smaller than
the distance to the fixed-gate ball boundary. A first-exit argument
then keeps the trajectory inside for all time. The risk gap and output
error tend to zero, proving attraction to $F_*$. The network need not
converge to one unique parameter realization.
\end{proof}
The same construction works with more neurons if their features remain
inactive at $x_*$ and a six-feature subcollection retains full rank.
One may positively balance each neuron without changing its output,
so this is compatible with the balanced parameterization used in the
experiments. It does not establish reachability from infinitesimal
isotropic Gaussian initialization.

\section{A fresh feature immediately exposes descent}
For any bias-free ReLU feature $h$ one has $h\in\mathcal B$.
Projection orthogonality shows that the residual functional on
$\mathcal B$ is proportional to evaluation at $x_*$:
\[
 \E_P[(f-F_*)h]=\frac{1-\rho}{7-\rho}\,h(x_*).             \tag{4}
\]
To identify the constant, use $h=f$, for which $f(x_*)=1$, and
$\E_P[(f-F_*)f]=\E_P(f-F_*)^2=2R(F_*)$.
Thus every incoming direction $v$ with $v\cdot x_*>0$ admits a fresh
correctly signed neuron giving descent. Appending parameters
$(a,w)=(t,tv)$ changes the output by $t^2\relu(v\cdot x)$ and changes
risk by exactly
\[
 -t^2\frac{1-\rho}{7-\rho}\relu(v\cdot x_*)
       +\frac{t^4}{2}\E_P\relu(v\cdot x)^2<0
\]
for sufficiently small $t>0$. This contrasts a stable local minimum
at fixed architecture with an immediate descent direction after adding
one neuron. A randomly signed isotropic fresh neuron has the needed
sign and activity with probability $1/4$; that observation alone is not
a recovery-time or full-training theorem.

\section{Numerical interpretation}
In the three-bit rare-seed experiment, the input distribution is written
as $(1+\rho x_1x_2x_3)/8$, with $\rho=.9$ and first-coordinate scale
$L=2$. Flipping suitable cube coordinates and, when necessary, the
output sign maps its low-probability dead vertices to the convention
above. Several width-1024 runs reach (1), with gradient norms around
$10^{-15}$ and risk agreeing with $1/122$ to displayed double precision.
Other runs hover near its nonsmooth boundary under a fixed time step.
These are numerical observations, not verified trajectories.

The proposed shortcut-escape criterion counted seeds with both readout
signs in a frozen residual field. This is insufficient to guarantee full
network recovery: the nonlinear evolution can instead enter the open
basin above. Counting all four parity-ray orientations is sufficient in
the completed sample of runs but not necessary; some runs generate
missing directions during training. Neither empirical statement has
been proved as a random-initialization theorem. Reported recovery should
require positive margins or near-zero square risk, not a floating-point
sign at an output mathematically equal to zero.
\end{document}
